\documentclass{article}

\textwidth=6.5in
\textheight=8.5in
\voffset=-54pt
\oddsidemargin=0in
\evensidemargin=0in
\setlength{\parskip}{8pt}
\setlength{\parindent}{0pt}

\usepackage{workbook}

\usepackage[sols]{handout}

\begin{document}

\begin{center}
  \large\textsc{Problem Set 1 - Relations, Equations, and Graphs}
\end{center}

\begin{enumerate}
\item 
\begin{grading}
    
\end{grading}

\begin{problem}
    Solve each equation for $y$ and state whether or not $y$ a function of $x$.
\end{problem}

\begin{enumerate}

    \item
    \begin{problem}
        $\ln(xy)=10$
    \end{problem}
    \pspace{\vfill}

    \begin{solution}
        To solve for $y$, we need to first get $y$ out of the logarithm by exponentiating both sides.
        \begin{align*}
            \ln(xy) &= 10\\
            e^{\ln(xy)} &= e^{10} \\
            xy &= e^{10}\\
            \intertext{Now we isolate $y$.}
            y &= \frac{e^{10}}{x}
        \end{align*}
        Since every value of $x$ corresponds to exactly one value of $y$, $y$ is a function of $x$ in this relationship.
    \end{solution}

    \item 
    \begin{problem}
        $x^2+y^2=25$
    \end{problem}
    \pspace{\vfill}

    \begin{solution}
        We can solve for $y$ by first isolating $y^2$.
        \begin{align*}
            x^2+y^2 &= 25 \\
            y^2 &= 25-x^2 \\
            \intertext{The solutions to this are}
            y=\pm \sqrt{25-x^2}
        \end{align*}
    One value of $x$ can correspond to more than one value of $y$. For example, $x=3$ corresponds to $y=\pm4$. So $y$ is not a function of $x$. 
    \end{solution}

    \item
    \begin{problem}
        $x^2y+4xy-1=0$
    \end{problem}
    \pspace{\vfill}

    \begin{solution}
    We want to isolate $y$.
    \begin{align*}
        x^2y+4xy-1&= 0 \\
        x^2y + 4xy &= 1 \\
        (x^2+4x)y &= 1 \\
        y &= \frac{1}{x^2+4x}
    \end{align*}
    Each value of $x$ corresponds to exactly one value of $y$, so $y$ is a function of $x$.
    \end{solution}

    \item 
    \begin{problem}
        $e^y - 2x = 1$
    \end{problem}
    \pspace{\vfill}

    \begin{solution}
        We can start by isolating $e^y$.
        \begin{align*}
            e^y &= 1+2x \\
            \intertext{Now we can take logs of both sides.}
            \ln(e^y) &= \ln(1+2x) \\
            y &= \ln(1+2x)
        \end{align*}
        Each value of $x$ corresponds to exactly one value of $y$, so $y$ is a function of $x$.
    \end{solution}

\end{enumerate}

\newpage

\item
\begin{problem}
In this problem, we will review how to find the distance between two points on a coordinate plane.
\end{problem}

\begin{enumerate}
    \item 
    \begin{problem}
    Shown in the diagram are the points $(-1, 2)$ and $(3, -1)$. A right triangle has been drawn with the segment between the two points as the hypotenuse. Determine the distance between the two points by finding the side lengths of the triangle.

    \begin{tikzpicture}
        \begin{axis}[
        xlabel=$x$, ylabel=$y$,
        xmin=-3.5, xmax=3.5, ymin=-3.5, ymax=3.5,
        xtick distance=1, ytick distance=1,
        xticklabels=\empty, yticklabels=\empty,
        axis equal image
        ]
            \addplot[draw=none]{x};
            \addplot[closed,mark size=2pt] coordinates {(-1,2) (3,-1)};
            \draw[very thick] (-1,2) -- (3, -1);
            \draw[thick, dashed] (-1,2) -- (3, 2);
            \draw[thick, dashed] (3,2) -- (3,-1);
            \draw[thick] (2.8,2)--(2.8,1.8)--(3,1.8);
            \solonly{
            \node[above] at (1,2) {$4$};
            \node[right] at (3,0.5) {$3$};
            \node[below left] at (1,0.5) {$c$};
            }
        \end{axis}
    \end{tikzpicture}
    \end{problem}

    \begin{solution}
        The length of the horizontal side of the right triangle is $4$, and the length of vertical side of the triangle is $3$. Using the Pythagorean theorem, we can determine the length of the hypotenuse $c$.
        \begin{align*}
        3^2 + 4^2 &= c^2 \\
        c^2 &= 25 \\
        c &= 5 \text{ (since $c$ needs to be positive)}
        \end{align*}
    The distance between the two points is 5 units.
    \end{solution}

    \item 
    \begin{problem}
    Here is a similar diagram, but we don't know the exact coordinates of the two points. Let's name the coordinates $(x_1, y_1)$ and $(x_2, y_2)$. Using the same process as the previous problem, come up with an expression for the distance between the two points. (Your answer should be in terms of the variables $x_1$, $x_2$, $y_1$, and $y_2$.)
    
    \begin{tikzpicture}
        \begin{axis}[
        xlabel=$x$, ylabel=$y$,
        xmin=-3.5, xmax=3.5, ymin=-3.5, ymax=3.5,
        xtick distance=1, ytick distance=1,
        xticklabels=\empty, yticklabels=\empty,
        axis equal image
        ]
            \addplot[draw=none]{x};
            \addplot[closed,mark size=2pt] coordinates {(-2.3,-1.4) (1.6,2.7)};
            \draw[very thick] (-2.3,-1.4) -- (1.6, 2.7);
            \node[below] at (-2.3,-1.4) {$(x_1, y_1)$};
            \node[right] at (1.6,2.7) {$(x_2,y_2)$};
            \draw[thick, dashed] (-2.3,-1.4) -- (1.6, -1.4);
            \draw[thick, dashed] (1.6,-1.4)--(1.6,2.7);
            \draw[thick] (1.4,-1.4)--(1.4,-1.2)--(1.6,-1.2);
            \solonly{
            \node[below] at (-0.35,-1.4) {$x_2-x_1$};
            \node[right] at (1.6, 0.65) {$y_2-y_1$};
            \node[above left] at (-0.35, 0.65) {$d$};
            }
        \end{axis}
    \end{tikzpicture}
    \end{problem}

    \begin{solution}
        The length of the horizontal side of the triangle is $x_2-x_1$, and the length of the vertical side of the triangle is $y_2-y_1$.
        We can use the Pythagorean theorem to solve for the length of the hypotenuse $d$.
        \begin{align*}
            d^2 &= (x_2-x_1)^2 + (y_2-y_1)^2  \\
            d &= \sqrt{(x_2-x_1)^2 + (y_2-y_1)^2} \text{ (since $d$ must be positive)}
        \end{align*}
    \end{solution}

    \item 
    \begin{problem}
        Using variables instead of specific values gives us a formula we can use to find the distance between \textbf{any} two points. We just need to plug in values for the variables. Verify that your answer to part (b) is correct by plugging in the points from part (a). You should get the same answer as before.
    \end{problem}
    \pspace{\vfill}

    \begin{solution}
        Let's say $(x_1, y_1)$ is $(-1, 2)$ and $(x_2, y_2)$ is $(3, -1)$. Plugging the values of $x_1$, $x_2$, $y_1$, and $y_2$ into the formula we found gives us
        \begin{align*}
            d &= \sqrt{(3-(-1))^2+(-1-2)^2} \\
            &= \sqrt{4^2+(-3)^2} \\
            &= \sqrt{25}\\
            &= 5
        \end{align*}
        This is the same as the answer we found in part (a).
    \end{solution}
\end{enumerate}

\pspace{\newpage}

\item
\begin{problem}
In this problem, we will investigate equations of circles that are \textbf{not} centered at the origin. Consider the circle whose center is $(2, 1)$ and whose radius is 1. You may wish to draw it on the axes provided.
\end{problem}

\begin{enumerate}
    \item
    \begin{problem}
        Determine the coordinates of four points that lie on the circle. Explain how you know.
    \end{problem}
    
     \begin{tikzpicture}
            \begin{axis}[
            xmin=-0.5, xmax=3.5, ymin=-0.5, ymax=3.5,
            axis equal image,
            xticklabels=\empty,
            yticklabels=\empty,
            xtick distance=1,
            ytick distance=1,
            xlabel=$x$, ylabel=$y$,
            width=2.2in
            ]
            \addplot[draw=none]{x};
            \solonly{\addplot[domain=0:360, very thick] ({2+cos(x)}, {1+sin(x)});}
            \addplot[closed] coordinates {(2,0) (1,1) (3,1) (2,2)};
            \end{axis}
        \end{tikzpicture}

        \begin{solution}
        All points on the circle must be a distance of 1 unit away from the center $(2,1)$. Four points that we know are 1 unit away from $(2,1)$ are $(2,0)$, $(1,1)$, $(3,1)$, and $(2,2)$.
    \end{solution}


    \item 
    \begin{problem}
         Is the point $(\frac{3}{2},\frac{1}{2})$ on the circle? Justify your answer using the distance formula.
    \end{problem}
    \pspace{\vfill}

    \begin{solution}
        To be on the circle, the point must be 1 unit away from the center. We can use the distance formula to find the distance between $(\frac{3}{2},\frac{1}{2})$ and $(2,1)$. 
        \begin{align*}
            d &= \sqrt{(2-\tfrac{3}{2})^2 + (1-\tfrac{1}{2})^2} \\
            &= \sqrt{(\tfrac{1}{2})^2 + (\tfrac{1}{2})^2} \\
            &= \sqrt{\tfrac{1}{4}+\tfrac{1}{4}}\\
            &= \sqrt{\tfrac{1}{2}}
        \end{align*}
        Since the distance from $(\frac{3}{2},\frac{1}{2})$ to the center of the circle is $\sqrt{1/2}$ and not 1, the point is not on the circle.
    \end{solution}

    \item 
    \begin{problem}
    Let $(x, y)$ be an arbitrary point on the circle. (It's somewhere on the circle, but we don't want to use specific values.) Write down an expression in terms of $x$ and $y$ that gives the distance between $(x, y)$ and the center of the circle $(2, 1)$.
    \end{problem}
    \pspace{\vfill}

    \begin{solution}
    We can use the distance formula with the points $(x, y)$ and $(2, 1)$ to get the expression $\sqrt{(x-2)^2+(y-1)^2}$.
    \end{solution}

    \item 
    \begin{problem}
    Based on your work in the previous parts, explain why the equation $(x-2)^2+(y-1)^2=1$ produces a circle centered at $(2, 1)$ with radius 1.
    \end{problem}
    \pspace{\vfill}

    \begin{solution}
        In part c, we determined that the distance between a point $(x, y)$ and $(2, 1)$ is $\sqrt{(x-2)^2+(y-1)^2}$. The circle centered at $(2, 1)$ with radius 1 consists of all points whose distance from $(2, 1)$ is 1 unit, which are all the solutions to the equation $\sqrt{(x-2)^2+(y-1)^2}=1$. Squaring both sides gives us the equation $(x-2)^2+(y-1)^2=1$.

        The solutions to the equation $(x-2)^2+(y-1)^2=1$ are all the points that are 1 unit away from $(2, 1)$, and this is exactly the circle centered at $(2, 1)$ with radius 1.
    \end{solution}
    

\end{enumerate}

\pspace{\newpage}

\item 
\begin{problem}
The equation $(x-h)^2+(y-k)^2=r^2$ is called the standard form of the equation of a circle with center $(h, k)$ and radius $r$. Use this fact to answer the following parts.
\end{problem}

\begin{enumerate}
    \item
    \begin{problem}
        Sketch the graph of $x^2+y^2=9$. Clearly indicate the center and radius. Find the exact coordinates of any $x$- and $y$-intercepts.
    \end{problem}
    \pspace{\vfill}

    \begin{solution}
            The graph of $x^2+y^2=9$ is a circle centered at the origin with radius 3.
    \begin{center}
        \begin{tikzpicture}
            \begin{axis}[xlabel=$x$, ylabel=$y$, axis equal image, xtick distance=1, ytick distance=1, xmin=-4.5, xmax=4.5, ymin=-4.5, ymax=4.5]
                \addplot[domain=0:360] ({3*cos(x)},{3*sin(x)});
            \end{axis}
        \end{tikzpicture}
    \end{center}
    The $x$-intercepts are $(-3,0)$ and $(3,0)$. The $y$-intercepts at $(0,-3)$ and $(0,3)$. We can see these from the graph, but we can also find them algebraically by setting $y=0$ and $x=0$. If $y=0$, then $x^2=9$, and the solutions are $x=\pm3$. If $x=0$, then $y^2=9$, and the solutions are $y\pm3$.
    \end{solution}

    \item
    \begin{enumerate}
    \item
    \begin{problem}
        Sketch the graph of $(x+3)^2+(y-6)^2=25$. Clearly indicate the center and radius. Find the exact coordinates of any $x$- and $y$-intercepts.
    \end{problem}
    \pspace{\vfill}

        \begin{solution}
            The graph of $(x+3)^2+(y-6)^2=25$ is a circle centered at $(-3,6)$ with radius 5.
    \begin{center}
        \begin{tikzpicture}
            \begin{axis}[xlabel=$x$, ylabel=$y$, axis equal image, extra x ticks={-8,-3,2}, extra y ticks={1,6,11}, xtick distance=1, ytick distance=1, xticklabels=\empty, yticklabels=\empty, xmin=-12.5, xmax=12.5, ymin=-12.5, ymax=12.5]
                \addplot[domain=0:360] ({5*cos(x)-3},{5*sin(x)+6});
                \addplot[closed] coordinates {(-3,6) (-8,6) (2,6) (-3,1) (-3,11)};
                % \node[above] at (-3,6) {$(-3,6)$};
                % \node[right] at (2,6) {$(2,6)$};
                \end{axis}
        \end{tikzpicture}
    \end{center}
    From the graph, we can see that there are no $x$-intercepts. We can verify this algebraically. If $y=0$,
    \begin{align*}
        (x+3)^2 + (0-6)^2 &= 25\\
        (x+3)^2 + 36 &= 25 \\
        (x+3)^2 &= -11\\
        \text{no real solutions}
    \end{align*}

    If $x=0$,
    \begin{align*}
        (0+3)^2 + (y-6)^2 &= 25 \\
        9 + (y-6)^2 &= 25\\
        (y-6)^2 &= 16 \\
        y-6 &= \pm4 \\
        y &= 6\pm4
    \end{align*}
    So there are two $y$-intercepts, $(0, 2)$ and $(0, 10)$, which matches the graph.
    \end{solution}

        \item
        \begin{problem}
        Find the $y$-coordinates of the points on the circle that have an $x$-coordinate of $-6$.
        \end{problem}
    \pspace{\vfill}

    \begin{solution}
        We can do this algebraically. If $x=-6$,
        \begin{align*}
            (-6+3)^2 + (y-6)^2 &= 25 \\
            9 + (y-6)^2 &= 25
        \end{align*}
        At this point, we can see that this equation is the same as an equation we solved in the previous part, and that the solutions were $y=6\pm4$. So the points on the graph with an $x$-coordinate of $-6$ are $(-6, 2)$ and $(-6, 10)$. This makes sense looking at the graph, and also is a result of the symmetry present. The $y$-intercepts we found before were 3 units to the left of the center, and the points we just found are 3 units to the right of the center.
    \end{solution}
    \end{enumerate}

\end{enumerate}

\pspace{\newpage}

\item
\begin{problem}
In this problem, we will sketch the graph of $y=x^3+3x^2$.
\end{problem}

\begin{enumerate}
    \item 
    \begin{problem}
        Find the $x$- and $y$-intercepts of the graph.
    \end{problem}
    \pspace{\vfill}

    \begin{solution}
        If $x=0$, then $y=0$, so the only $y$-intercept is $(0,0)$. 

        If $y=0$, then we have
        \begin{align*}
            0&= x^3+3x^2 \\
            &= x^2(x+3) \\
            x&=0, -3
        \end{align*}
        So the $x$-intercepts are $(0,0)$ and $(-3,0)$.
    \end{solution}
    

    \item 
    \begin{problem}
        Find the intervals on which $f(x)=x^3+3x^2$ is increasing and decreasing.
    \end{problem}
    \pspace{\stretch{2}}

    \begin{solution}
        First we need the derivative. Using the Power Rule, $f'(x)=3x^2+6x$. Now we want to analyze the sign of the derivative. If we factor the right side, we have $f'(x)=3x(x+2)$, so $f'(x)=0$ when $x=0,-2$. We can make a sign-chart by checking the sign of $f'(x)$ or using our knowledge of the graph.

        \begin{center}
            \begin{tikzpicture}
    % Draw the number line
    \draw[<->] (-4, 0) -- (2, 0) node[right] {$x$};

    % Points for critical values
    \foreach \x in {-2, 0} {
        \draw (\x,0.1) -- (\x,-0.1); % Draw vertical tick marks
    }

    % Labels for critical points
    \node at (-2, -0.3) {\small $-2$};
    \node at (0, -0.3) {\small $0$};

    % Mark intervals and indicate sign
    \node at (-3, 0.3) {\small $+$};

    \node at (-1, 0.3) {\small $-$};

    \node at (1, 0.3) {\small $+$};

    % Indicate zero points
    \node at (-2, 0.3) {\small $0$};
    \node at (0, 0.3) {\small $0$};
    
    % Label the function
    \node at (-4.3,0.3) {$f'(x)$};
\end{tikzpicture}
        \end{center}
        Since $f'(x)$ is positive on the intervals $(-\infty, -2)$ and $(0, \infty)$, $f(x)$ is increasing on the intervals $(-\infty, 2)$ and $(0,\infty)$. Since $f'(x)$ is negative on the interval $(-2,0)$, $f'(x)$ is decreasing on the interval $(-2,0)$.
    \end{solution}

    \item
    \begin{problem}
        Sketch the graph. Clearly indicate the coordinates of all intercepts and extrema.
    \end{problem}
    \pspace{\nstretch{2}}

    \begin{solution}
    Since this is a cubic function with a positive leading coefficient, the ends of the graph should go down and to the left and up and to the right. Coming from the left, the function increases, passing through the $x$-axis at $x=-3$, until it reaches a local maximum at $(-2,4)$. The function then decreases until it reaches a local mininum at $(0,0)$, then increases. 
        \begin{center}
            \begin{tikzpicture}
                \begin{axis}[xlabel=$x$, ylabel=$y$, extra x ticks={-3, -2}, extra y ticks={4}, xtick=\empty, ytick=\empty]
                    \addplot[domain=-3.5:1.5] {x^3+3*x^2};
                    \addplot[closed] coordinates {(-3,0) (0,0) (-2,4)};
                \end{axis}
            \end{tikzpicture}
        \end{center}
    \end{solution}

    \item
    \begin{problem}
        Is it possible to rewrite $y=x^3+3x^2$ into the form $x=g(y)$? Briefly explain why or why not.
    \end{problem}
    \pspace{\vfill}

    \begin{solution}
        No, it's not possible. From the graph, we can see that a single $y$-value can correspond to multiple $x$-values, meaning that it's not possible for us to express $x$ as a function of $y$.
    \end{solution}
    
\end{enumerate}

\item 
  \note{This is a quick warm-up to help get students comfortable with
    differentiating in both notations.
    
    Matt note - at the end of this problem, I'll do a quick question asking them what the relationship is between 1a and 1b.}

  \begin{enumerate}
  \item
    \begin{problem}[\vfill]
      If $g(t) = [f(t)]^3$, express $g'(t)$ in terms of $f(t)$ and $f'(t)$.
    \end{problem}

    \begin{solution}
      By the Chain Rule, $\displaystyle g'(t) = \frac{d}{dt}
      \left(\fcolorbox{red}{white}{$\displaystyle
          \fcolorbox{blue}{white}{$\displaystyle f(t)$}^ { 3} $}\right) =
      3 [f(t)]^2 \cdot f'(t)$.
    \end{solution}

  \item
    \begin{problem}[\vfill]
      If $y=x^3$ and $x$ is a function of $t$, express $\frac{dy}{dt}$
      in terms of $x$ and $\frac{dx}{dt}$. 
    \end{problem}

    \begin{solution}
      This is actually the same question as
      {{{{\#1}}(a)}}, just in different notation.
      Since $x$ is an unknown function of $t$, we can think of $x$ as
      $f(t)$, even though we don't know what function $f$ is specifically. The notation $\frac{dx}{dt}$ represents the same derivative as $f'(t)$.

      \begin{align*}
          \frac{dy}{dt} &= \frac{d}{dt} x^3 \\
          &= \frac{d}{dt} [f(t)]^3 \\
          &= 3[f(t)]^2 \cdot f'(t)\\
          \frac{dy}{dt}&= 3x^2\frac{dx}{dt}
      \end{align*}
      
    \end{solution}

  \end{enumerate}

\end{enumerate}

\psreflection{
  \emph{Reminders:}
  \begin{itemize}[topsep=0pt]
  \item Section for workshop ASAP if you haven't already! Our first workshop is on \textbf{Tuesday, February 4th}.
  \item Complete Math Ma Skills Check Practice 1 and turn it in by \textbf{Thursday, January 30th at 5:30pm.}
  \item The Math Ma Skills Check
  is \textbf{Thursday, February 6th 5:30-6:30pm}.  If you have an approved conflict (please see syllabus for what conflicts will be approved), you must fill out the Out of Sequence Exam form (found on Canvas) by Friday, January 31st at 7pm to arrange an earlier time to take the Skills Check.
  \end{itemize}

  Example 17.9 on pg. 551 of Gottlieb is recommended reading for next class.
}

\end{document}
